\documentclass[11pt,letterpaper]{article}
\usepackage[T1]{fontenc}
\usepackage{lmodern}
\usepackage[margin=1in]{geometry}
\usepackage{microtype}
\usepackage{amsmath,amssymb,amsthm,mathtools}
\usepackage{enumitem}
\usepackage{xcolor}
\usepackage{hyperref}
\hypersetup{colorlinks=true,linkcolor=blue!45!black,citecolor=blue!45!black,
  urlcolor=blue!45!black,pdftitle={Golden-ratio growth of Conway's subprime closure},
  pdfauthor={Romain Popescu},
  pdfsubject={A proof for the cardinality-ratio conjecture of Caragiu, Vicol, and Zaki}}
\setlist[enumerate]{leftmargin=*,itemsep=0.65em,topsep=0.6em}
\numberwithin{equation}{section}
\newtheorem{theorem}{Theorem}[section]
\newtheorem{lemma}[theorem]{Lemma}
\newtheorem{proposition}[theorem]{Proposition}
\theoremstyle{remark}
\newtheorem{remark}[theorem]{Remark}
\DeclareMathOperator{\lpf}{lpf}
\DeclareMathOperator{\meas}{meas}
\newcommand{\N}{\mathbb N}
\newcommand{\Z}{\mathbb Z}
\newcommand{\R}{\mathbb R}
\newcommand{\Pp}{\mathbb P}
\newcommand{\ph}{\varphi}
\newcommand{\tot}{\phi_{\!E}}
\newcommand{\ind}{\mathbf 1}
\newcommand{\ee}[1]{e\!\left(#1\right)}
\newcommand{\floor}[1]{\lfloor #1\rfloor}
\newcommand{\ceil}[1]{\lceil #1\rceil}
\newcommand{\HH}[2]{H_{#1}(#2)}
\allowdisplaybreaks[2]
\title{Golden-ratio growth of\\Conway's subprime closure}
\author{Romain Popescu}
\date{12 September 2026}

\begin{document}

\maketitle


\begin{abstract}
Let $s(m)$ be the Conway subprime function  \cite{OEIS}, \cite{Roberts}, define the binary operation on the natural numbers
$x \circ y= s(x + y)$, and denote by $C_n$, $n \ge 0$,
the sequence of subsets of natural numbers defined by $C_0 = \{1\}$,
and $C_{n+1} = C_n \cup (C_n \circ C_n)$. We prove the conjecture \cite{CVZ} by
Caragiu, Vicol and Zaki that
\[
\lim_{n\to \infty} \frac{ | C_{n+1}| }{ | C_n |}= \frac{1+\sqrt{5}}{2}.
\]
The underlying mathematical proof in this paper was constructed with some algorithmic assistance from GPT-6 Astra \cite{openai2026astra}.
The main theorem has been formally verified in the Lean proof assistant, taking as analytic input the Siegel--Walfisz theorem and a Vinogradov-type minor-arc bound; the formalization is available at \cite{Github}.
\end{abstract}

\noindent\textbf{2020 Mathematics Subject Classification.}
11N05, 11P32, 11P55, 11B39.

\noindent\textbf{Key words and phrases.}
Conway subprime function, golden ratio, Goldbach-type theorems, Hardy--Littlewood circle method, Lean.

\section{Introduction}

For a positive integer $m$, the Conway's subprime function \cite{OEIS}, \cite{Roberts} is defined as follows
\[
 s(m)=
 \begin{cases}
  1,&m=1,\\
  m,&m\text{ is prime},\\
  m/\lpf(m),&m\text{ is composite},
 \end{cases}
\]
where $\lpf(m)$ is the least prime factor of $m$. 

The Conway subprime function is used in the construction of  subprime Fibonacci sequences, i.e. integer
sequences $\{x_k\}$, $k\ge 0$,  satisfying a recursion $x_k = s (x_{k-1}+x_{k-2})$,
which have been studied (see \cite{GuyKhovanovaSalazar2014})
for periodicity, in analogy with the Collatz problem.


In  \cite{CVZ}, Caragiu, Vicol, and Zaki  use the Conway subprime function to construct the following sets.
Starting with $1$, retain all previously generated numbers and apply the operation $(a,b)\mapsto s(a+b)$
to every pair available at the preceding stage:
\begin{equation}\label{eq:closure}
 C_0=\{1\},\qquad
 C_{n+1}=C_n \cup (C_n \circ C_n) = C_n\cup\{s(a+b):a,b\in C_n\}.
\end{equation}

The main result of \cite{CVZ} is that

\begin{equation}\label{eq:exhaustion}
 \bigcup_{n\ge0}C_n=\N,
\end{equation}
%
The golden ratio $\ph=\frac{1+\sqrt5}{2}$ makes an unexpected appearance in their paper, since
Caragiu, Vicol, and Zaki  conjecture in \cite{CVZ} that

\begin{equation}\label{eq:conjecture}
 \lim_{n\to\infty}\frac{|C_{n+1}|}{|C_n|}
 =\ph.
\end{equation}



Write $M_n=\max C_n$. For $n\ge1$, let $R_n$ be the first prime absent
from $C_n$, and let $Q_n$ be the prime immediately preceding $R_n$.
These are well-defined: $C_n$ is finite and contains $2$. Thus every prime
at most $Q_n$ belongs to $C_n$, while $R_n\notin C_n$ and $Q_n\le M_n$.
We prove the following theorem, which is a more precise version of the aforementioned conjecture 3 in \cite{CVZ}:

\begin{theorem}\label{thm:main}
There is a constant $c>0$ such that
\begin{equation}\label{eq:main-scales}
 M_n\sim Q_n\sim c\ph^n,\qquad |C_n|\sim c\ph^{n-1}.
\end{equation}
Consequently,
\begin{equation}\label{eq:main-limits}
 \lim_{n\to\infty}\frac{|C_{n+1}|}{|C_n|}=\ph,
 \qquad
 \lim_{n\to\infty}\frac{|C_n|}{M_n}=\frac1\ph.
\end{equation}
\end{theorem}

The proof of the  theorem uses the cited exhaustion result in \cite{CVZ},
an almost-all theorem for Goldbach representations with nearly equal primes,
and two standard estimates for exponential sums over primes. The particular
interval-restricted representation lemma needed here is proved in
Appendix~\ref{app:binary}. The proof was constructed with algorithmic assistance from OpenAI Astra.
A Lean formalization of the main theorem, assuming the Siegel--Walfisz theorem and a Vinogradov-type minor-arc bound, is available at \cite{Github}.

\section{Heuristics and Proof overview}

\begin{enumerate}
\item \textbf{Parity gives the Fibonacci upper bound.}
If a sum is composite, division by its least prime factor prevents the
output from exceeding the preceding maximum. A new maximum must therefore
be prime. An odd prime is a sum of an odd and an even input, and the even
input cannot exceed the maximum from one generation earlier. This gives
$M_{n+1}\le M_n+M_{n-1}$.

\item \textbf{Complete prime prefixes reproduce at the Fibonacci scale.}
When every prime up to $X$ is present, almost-all Goldbach representations
produce all but $O_A(X/\log^A X)$ integers up to $X$ one generation later.
These integers, paired with primes near a larger cutoff, extend the complete
prime prefix by an almost Fibonacci step. Summable relative losses then give
positive constants $c,d$ with $M_n\sim c\ph^n$ and $Q_n\sim d\ph^n$.

\item \textbf{A gap between these constants produces extra primes.}
If $d<c$, an earlier maximum supplies a generated prime $t$ beyond the
predicted complete prefix at a suitable scale. The equation
$2p+q=2u-t$, with $p$ and $q$ in specified complete prime intervals,
generates $u$ through the two operations
$(t+q)/2$ and $p+(t+q)/2$. An almost-all representation estimate therefore
produces almost every prime in an interval beyond the predicted prefix.

\item \textbf{Two small exceptional sets eliminate the gap.}
Take the first missing prime two generations later. It has many possible
decompositions $r=u+(r-u)$ with $u$ in the extra interval. Only a small
number of these $u$ are missing, and only a small number of their complementary
integers $r-u$ are missing. At least one pair is present, forcing $r$ to be
generated and contradicting its definition. Hence $d=c$.

\item \textbf{The preceding maximum determines the cardinality.}
Since $Q_{n-1}\sim M_{n-1}$, the Goldbach filling estimate makes
$C_n$ almost complete below $M_{n-1}$. Every element above that threshold
is prime, and those primes contribute only $o(M_{n-1})$ elements.
Thus $|C_n|\sim M_{n-1}$, yielding both the golden-ratio growth and
the limiting density $1/\ph$.
\end{enumerate}

\section{Analytic preliminaries}

All logarithms are natural. We will denote cardinality of a set $X$ either by $|X|$ or by $\# X$. \ 
The notations $f=O_A(g)$ and $f\ll_A g$ mean $|f|\le C_Ag$ for all sufficiently large values of the relevant
parameter. We will write $f\asymp g$ when both $f\ll g$ and $g\ll f$ hold, and $f\sim g$
when $f/g\to1$. 

Let $\Pp$ denote the set of primes and $\pi(X)$
the number of primes at most $X$. Intervals used to count missing
elements are to be intersected with $\N$. Define
\[
 H_j(X)=\#\bigl(\{1,\ldots,\floor X\}\setminus C_j\bigr).
\]

\begin{lemma}[Prime counts]\label{lem:pnt}
As $X\to\infty$,
\begin{equation}\label{eq:short-primes}
 \#\left(\Pp\cap\left[X-\frac{X}{\log^2X},X\right]\right)
 \sim\frac{X}{\log^3X}.
\end{equation}
For fixed $0<\alpha<\beta$,
\begin{equation}\label{eq:proportional-primes}
 \#\bigl(\Pp\cap[\alpha X,\beta X]\bigr)
 \sim\frac{(\beta-\alpha)X}{\log X}.
\end{equation}
Also $\pi(X)\sim X/\log X$. The prime immediately preceding a large
real $X$ is $X+O(X/\log^2X)$, and the prime immediately following $X$
is $X+o(X)$.
\end{lemma}

\begin{proof}
The $q=1$, zero-frequency case of \cite[Lemma 5 (Major arc estimates), equation (26)]{LyallRice} or \cite[Proposition~24]{Tao},  
which both follow from the Siegel-Walfisz theorem on primes in arithmetic progressions, yields
for every fixed $D>0$, that
\[
 \sum_{m\le X}\Lambda(m)=X+O_D(X/\log^D X).
\]
The prime powers with exponent at least two have total weight
$O(\sqrt X\log^2X)$: there are at most $\log_2X$ possible exponents,
and for each exponent at most $\sqrt X$ bases, each of weight at most
$\log X$. Thus
\begin{equation}\label{eq:theta}
 \vartheta(X):=\sum_{p\le X}\log p
 =X+O_D(X/\log^D X).
\end{equation}
Subtract this formula at $X$ and $X-h$, where $h=X/\log^2X$, taking
$D=4$. It gives total prime weight $h+O(X/\log^4X)\sim h$ in
$(X-h,X]$. Every such prime has logarithm $\sim\log X$, uniformly,
which proves \eqref{eq:short-primes}; changing an endpoint adds at most
one prime. The same subtraction at $\beta X$ and $\alpha X$ proves
\eqref{eq:proportional-primes}. Partial summation gives
\[
 \pi(X)=\frac{\vartheta(X)}{\log X}
       +\int_2^X\frac{\vartheta(t)}{t\log^2t}\,dt
 =\frac{X}{\log X}+O(X/\log^2X).
\]
For the integral bound, split at $\sqrt X$ and use $\vartheta(t)\ll t$.
The preceding-prime claim follows from the nonemptiness of the interval
in \eqref{eq:short-primes}. For every fixed $\epsilon>0$,
\eqref{eq:proportional-primes} supplies a prime between $X$ and
$(1+\epsilon)X$ for all sufficiently large $X$. Letting $\epsilon$
be arbitrary proves the following-prime claim.
\end{proof}

\begin{lemma}[Goldbach filling]\label{lem:filling}
Fix $A>0$. If $C_j$ contains every prime at most $X$, then
\begin{equation}\label{eq:filling}
 H_{j+1}(X)\ll_A\frac{X}{\log^A X}.
\end{equation}
The implicit constant is independent of $j$.
\end{lemma}

\begin{proof}
The corollary to Theorem~1 in \cite{CL}, with exponent $5/8+1/8=3/4$, or equations (1.5), (1.6) and Theorem 1.2 in \cite{BakerHarman1998}
yield that all but
$O_A(Y/\log^A Y)$ integers $m\in[Y,2Y]$ have a representation
\begin{equation}\label{eq:near-goldbach}
 2m=p+q,\qquad p,q\in\Pp,\qquad |p-m|,|q-m|\le m^{3/4}.
\end{equation}
Here their parameter for the even number $2m$ is rescaled by a factor
of two, which does not change the order of the exceptional-set bound.

Consider $\sqrt X<m\le X-X^{3/4}$. Cover this range by dyadic
intervals $[X/2^{r+1},X/2^r]$ that meet it. Each lower endpoint $Y$
is at least $\sqrt X/2$, their sum is at most $X$, and therefore
$\log Y\ge \tfrac13\log X$ for large $X$. The total number of
exceptions to \eqref{eq:near-goldbach} is consequently at most
\[
 C_A\sum_r\frac{X/2^{r+1}}{\log^A(X/2^{r+1})}
 \le 3^AC_A\frac{X}{\log^A X}.
\]
For every nonexceptional $m$ in this range,
$p,q\le m+m^{3/4}\le X$. Hence $p,q\in C_j$, and, as $m\ge2$,
the number $2m$ is composite with least prime factor $2$. Thus
$s(p+q)=m\in C_{j+1}$. The omitted bottom and top ranges contain
$O(\sqrt X+X^{3/4}+1)$ integers. This is $O_A(X/\log^A X)$,
proving the lemma.
\end{proof}

\begin{lemma}[Two primes in specified intervals]\label{lem:binary}
Fix $0<\alpha<\beta$, $0<\gamma<\delta$, $\eta>0$, $B>0$, and $A>0$.
Put
\[
 I_X=[\alpha X,\beta X],\qquad J_X=[\gamma X,\delta X],
 \qquad
 \mathcal J_X(N)=\meas\{v\in I_X:N-2v\in J_X\}.
\]
Among the odd integers $1\le N\le BX$ with
$\mathcal J_X(N)\ge\eta X$, at most $O(X/\log^A X)$ fail to have
a representation
\begin{equation}\label{eq:binary-rep}
 N=2p+q,\qquad p\in\Pp\cap I_X,\quad q\in\Pp\cap J_X.
\end{equation}
The implicit constant may depend on all the displayed fixed parameters.
\end{lemma}

The proof of this lemma is given in Appendix~\ref{app:binary}. It uses no assertion
about every individual odd target; the exceptional set is essential.

\section{Fibonacci growth and the two asymptotic constants}

\begin{lemma}[Structure and upper bound]\label{lem:structure}
For $n\ge1$, the maximum $M_n$ is prime and
\begin{equation}\label{eq:structure}
 C_n\subseteq\{1,\ldots,M_{n-1}\}
       \cup\bigl(\Pp\cap[1,M_n]\bigr).
\end{equation}
For $n\ge1$,
\begin{equation}\label{eq:fib-upper}
 M_{n+1}\le M_n+M_{n-1},\qquad M_n\le F_{n+2},
\end{equation}
where $F_0=0$, $F_1=1$, and $F_{r+2}=F_{r+1}+F_r$.
\end{lemma}

\begin{proof}
If $a,b\le M_{n-1}$ and $a+b$ is composite, then
\[
 s(a+b)=\frac{a+b}{\lpf(a+b)}\le\frac{a+b}{2}\le M_{n-1}.
\]
If $a+b$ is prime, the output is prime. Previously retained elements
are also at most $M_{n-1}$, proving \eqref{eq:structure}. Since
$M_1=2$, a maximum that changes can change only to a prime, and a
maximum that does not change remains prime. Moreover $M_2=3$, so
$M_n$ is odd for $n\ge2$.

For $n\ge2$, every even element of $C_n$ is at most $M_{n-1}$:
this follows from \eqref{eq:structure} for composites, and from
$2\le M_{n-1}$ for the only even prime. If $M_{n+1}>M_n$, its
generating sum is an odd prime, so one input is even and at most
$M_{n-1}$, and the other is at most $M_n$. If the maximum does not
increase, the same bound is immediate. The case $n=1$ is
$M_2=3=M_1+M_0$. Induction from $M_0=1$, $M_1=2$ now proves
the Fibonacci bound. In particular $M_n\ll\ph^n$, for example by
the formula $F_r=(\ph^r-(-\ph)^{-r})/\sqrt5$.
\end{proof}

\begin{lemma}[Extension of a complete prime prefix]\label{lem:extension}
Fix $B\ge1$. Suppose $Y\le X\le BY$, $C_n$ contains every prime
at most $X$, and $C_{n-1}$ contains every prime at most $Y$.
For sufficiently large $Y$, depending only on $B$, the set $C_{n+1}$
contains every prime at most
\begin{equation}\label{eq:extension}
 X+Y-\frac{X}{\log^2X}.
\end{equation}
\end{lemma}

\begin{proof}
Set $h=X/\log^2X$. For large $Y$, $h<Y$. Lemma~\ref{lem:filling}
gives $H_n(Y)\ll Y/\log^4Y$. Let $r$ be a prime with
$X<r\le X+Y-h$. For each prime $p\in[X-h,X]$, the integer
$a=r-p$ is positive and satisfies
\[
 a\le(X+Y-h)-(X-h)=Y.
\]
Distinct choices of $p$ give distinct $a$. By Lemma~\ref{lem:pnt},
the number of candidate primes is asymptotic to $X/\log^3X$, which
exceeds $H_n(Y)$ for sufficiently large $Y$: indeed $X\asymp_B Y$
and $\log X\sim\log Y$. Some candidate therefore has $a\in C_n$.
Also $p\in C_n$, and $s(p+a)=r$ because $r$ is prime. The primes
at most $X$ are retained, completing the proof.
\end{proof}

\begin{proposition}\label{prop:rough-scales}
We have $Q_n\asymp M_n\asymp\ph^n$.
\end{proposition}

\begin{proof}
Choose a fixed real $L$ large enough for Lemma~\ref{lem:extension}
with $B=2$, and such that $1/\log^2L\le1/4$.
By \eqref{eq:exhaustion} and nesting there is an $N$ with
$\{1,\ldots,\ceil L\}\subseteq C_N$. Define
\begin{equation}\label{eq:L-recursion}
 L_0=L_1=L,\qquad
 L_{k+1}=L_k+L_{k-1}-\epsilon_kL_k,\qquad
 \epsilon_k=\frac1{\log^2L_k}\quad(k\ge1).
\end{equation}
Inductively $L_{k-1}\le L_k\le2L_{k-1}$ and $L_k\ge L$.
Indeed, under these inequalities, $\epsilon_kL_k\le L_{k-1}/2$,
and therefore
\begin{equation}\label{eq:L-growth}
 \frac54L_k\le L_k+\frac12L_{k-1}
 \le L_{k+1}\le L_k+L_{k-1}\le2L_k.
\end{equation}
The base case is $L_0=L_1$. Lemma~\ref{lem:extension}, applied
successively, shows that every prime up to $L_k$ lies in $C_{N+k}$.
The two initial cases follow from the choice of $N$ and retention.

By \eqref{eq:L-growth}, $L_k\ge L(5/4)^{k-1}$ for $k\ge1$,
so $\sum_{k\ge1}\epsilon_k<\infty$. Set $u_k=L_k/\ph^k$.
Since $\ph^{-1}+\ph^{-2}=1$, \eqref{eq:L-recursion} implies
\[
 u_{k+1}\ge(1-\epsilon_k)
 \left(\frac{u_k}{\ph}+\frac{u_{k-1}}{\ph^2}\right).
\]
For $m_k=\min(u_{k-1},u_k)$ this gives
$m_{k+1}\ge(1-\epsilon_k)m_k$. Thus
\[
 m_k\ge m_1\prod_{j=1}^{k-1}(1-\epsilon_j)
 \ge m_1\exp\left(-2\sum_{j\ge1}\epsilon_j\right)>0,
\]
where $\log(1-z)\ge-2z$ for $0\le z\le1/2$.
Hence $L_k\gg\ph^k$. The largest prime at most $L_k$ is
$\sim L_k$ by Lemma~\ref{lem:pnt}, and is at most $Q_{N+k}$.
Consequently $Q_n\gg\ph^n$. Together with $Q_n\le M_n\ll\ph^n$,
this proves the proposition.
\end{proof}

\begin{lemma}[A contracting recurrence]\label{lem:contraction}
Suppose a bounded real sequence $(x_n)$ satisfies
$x_n+\rho x_{n-1}\to t$ for some $0<\rho<1$.
Then $x_n\to t/(1+\rho)$.
\end{lemma}

\begin{proof}
Put $y_n=x_n-t/(1+\rho)$. Then $y_n=-\rho y_{n-1}+o(1)$.
The finite number $L=\limsup_n|y_n|$ satisfies $L\le\rho L$,
so $L=0$.
\end{proof}

\begin{proposition}\label{prop:constants}
There are constants $0<d\le c$ such that
\begin{equation}\label{eq:constants}
 M_n\sim c\ph^n,\qquad Q_n\sim d\ph^n.
\end{equation}
\end{proposition}

\begin{proof}
Define
\[
 T_n=\ph^{-n}\left(M_n+\frac{M_{n-1}}\ph\right).
\]
Using $\ph-\ph^{-1}=1$ and \eqref{eq:fib-upper}, we obtain
\[
 T_{n+1}-T_n
 =\ph^{-(n+1)}(M_{n+1}-M_n-M_{n-1})\le0.
\]
Proposition~\ref{prop:rough-scales} bounds $T_n$ away from zero,
so $T_n\to t>0$. Apply Lemma~\ref{lem:contraction} to
$x_n=M_n/\ph^n$ and $\rho=\ph^{-2}$. It gives
$M_n/\ph^n\to c=t/(1+\ph^{-2})>0$.

The sequence $Q_n$ is nondecreasing, and
Proposition~\ref{prop:rough-scales} bounds $Q_n/Q_{n-1}$ above by
a fixed constant. Apply Lemma~\ref{lem:extension} with
$X=Q_n$, $Y=Q_{n-1}$. All primes up to
$Z=Q_n+Q_{n-1}-Q_n/\log^2Q_n$ belong to $C_{n+1}$.
Since $Z\asymp Q_n$, Lemma~\ref{lem:pnt} shows that the largest
prime at most $Z$ is at least $Z-O(Q_n/\log^2Q_n)$. Therefore
\begin{equation}\label{eq:Q-lower-recursion}
 Q_{n+1}\ge Q_n+Q_{n-1}-O(Q_n/\log^2Q_n).
\end{equation}
Set
\[
 V_n=\ph^{-n}\left(Q_n+\frac{Q_{n-1}}\ph\right).
\]
It is bounded and bounded away from zero. Since
$Q_n\asymp\ph^n$ and $\log Q_n=n\log\ph+O(1)$,
\eqref{eq:Q-lower-recursion} gives, for fixed $D>0$ and all $n\ge n_0$,
\[
 V_{n+1}-V_n\ge-\frac{D}{n^2}.
\]
The sequence $V_n+D\sum_{j=n_0}^{n-1}j^{-2}$ is bounded and
nondecreasing, hence convergent. Since the added sums converge,
$V_n\to v>0$. Lemma~\ref{lem:contraction} now gives
$Q_n/\ph^n\to d=v/(1+\ph^{-2})>0$. Finally $Q_n\le M_n$
implies $d\le c$.
\end{proof}

\section{Approximate completeness of the prime prefix}

\begin{proposition}\label{prop:equal}
The constants in Proposition~\ref{prop:constants} are equal:
$d=c$. In particular $Q_n/M_n\to1$.
\end{proposition}

\begin{proof}
Assume $d<c$. Choose a fixed integer $k\ge0$ such that
\begin{equation}\label{eq:choose-T}
 d<T:=c\ph^{-k}\le\ph d.
\end{equation}
For example, $k=\ceil{\log(c/d)/\log\ph}-1$ has these properties.
Let
\[
 X=\ph^{n-1},\qquad t_n=M_{n-1-k},\qquad
 \Delta=T-d,\qquad \epsilon=\Delta/16.
\]
All these constants except $X$ and $t_n$ are now fixed.
For sufficiently large $n$, the integer $t_n$ is an odd prime in
$C_{n-1}$, by Lemma~\ref{lem:structure} and retention. Moreover,
\begin{equation}\label{eq:three-scales}
 t_n=TX+o(X),\qquad Q_{n-1}=dX+o(X),\qquad
 Q_n=\ph dX+o(X).
\end{equation}

\medskip\noindent\textit{Generating primes beyond the predicted prefix.}
Define
\begin{equation}\label{eq:three-intervals}
 \begin{aligned}
 I_X&=[(\ph d-2\epsilon)X,(\ph d-\epsilon)X],\\
 J_X&=[(2d-T+3\epsilon)X,(2d-T+9\epsilon)X],\\
 U_X&=[(\ph^2d+\epsilon)X,(\ph^2d+2\epsilon)X].
 \end{aligned}
\end{equation}
All endpoints are positive. In particular, $T\le\ph d<2d$ makes
the lower endpoint of $J_X$ positive, and
$\epsilon\le(\ph-1)d/16$ makes that of $I_X$ positive.
The upper endpoint of $I_X$ is below $Q_n$ for large $n$, and
the upper endpoint of $J_X$ is
\[
 (2d-T+9\epsilon)X=(d-7\epsilon)X<Q_{n-1}.
\]
It follows that every prime in $I_X$ belongs to $C_n$, and every
prime in $J_X$ belongs to $C_{n-1}$.

For an integer $u\in U_X$, consider the target $N=2u-t_n$.
It is odd, and it lies in $[1,BX]$ for a fixed sufficiently large
$B$: its leading lower coefficient is
$2\ph^2d+2\epsilon-T>0$, and its upper coefficient is fixed.
For every real $p\in I_X$, the value $q=N-2p$ satisfies
\begin{equation}\label{eq:q-margins}
 (2d-T+4\epsilon)X+o(X)
 \le q\le
 (2d-T+8\epsilon)X+o(X).
\end{equation}
To verify these bounds, use the lower endpoint of $U_X$ and the
upper endpoint of $I_X$ for the lower bound, and the opposite
endpoints for the upper bound, together with $\ph^2-\ph=1$.
The error is exactly $-(t_n-TX)$, independent of $u$ and $p$.
For large $n$ its absolute value is less than $\epsilon X/2$.
Thus every such $q$ lies in $J_X$, and
\[
 \meas\{p\in I_X:2u-t_n-2p\in J_X\}=|I_X|=\epsilon X.
\]

Lemma~\ref{lem:binary} applies with these fixed intervals and
$\eta=\epsilon$. Since $u\mapsto2u-t_n$ is injective, all but
$O_A(X/\log^A X)$ integers $u\in U_X$ have a representation
\begin{equation}\label{eq:bridge}
 2p+q=2u-t_n,\qquad p\in\Pp\cap I_X,\quad q\in\Pp\cap J_X.
\end{equation}
If such a $u$ is itself prime, both operations in
\begin{equation}\label{eq:two-operations}
 b=\frac{t_n+q}{2}=s(t_n+q)\in C_n,
 \qquad u=p+b=s(p+b)\in C_{n+1}
\end{equation}
are valid. Indeed $t_n$ and $q$ are odd primes in $C_{n-1}$,
so their sum is composite with least prime factor $2$; and
$p\in C_n$, while the final sum $u$ is prime. Consequently
\begin{equation}\label{eq:extra-primes}
 \#\bigl((\Pp\cap U_X)\setminus C_{n+1}\bigr)
 \ll_A X/\log^A X.
\end{equation}

%\medskip\noindent\textit{Contradicting the first missing prime.}
Let $r_n=R_{n+2}$, the first prime absent from $C_{n+2}$.
It is the prime immediately after $Q_{n+2}$. By
Lemma~\ref{lem:pnt} and Proposition~\ref{prop:constants},
\begin{equation}\label{eq:first-missing-scale}
 r_n=\ph^3dX+o(X).
\end{equation}
For a prime $u\in U_X$, put $a=r_n-u$. The identity
$\ph^3-\ph^2=\ph$ gives
\[
 (\ph d-2\epsilon)X+o(X)
 \le a\le(\ph d-\epsilon)X+o(X).
\]
The errors here are uniform in $u$. By \eqref{eq:three-scales}
and the fixed positive margins, $1\le a\le Q_n$ for large $n$.

Lemma~\ref{lem:pnt} supplies
\[
 \#(\Pp\cap U_X)\sim\frac{\epsilon X}{\log X}
\]
candidate primes. At most $O(X/\log^4X)$ of them are absent from
$C_{n+1}$, by \eqref{eq:extra-primes} with $A=4$. Among the
remaining candidates, at most $H_{n+1}(Q_n)$ have
$r_n-u\notin C_{n+1}$, since distinct $u$ give distinct differences.
Lemma~\ref{lem:filling} and $Q_n\asymp X$ give
\[
 H_{n+1}(Q_n)\ll\frac{Q_n}{\log^4Q_n}
 \ll\frac{X}{\log^4X}.
\]
The combined exclusions are $o(X/\log X)$, so at least one candidate
has both $u\in C_{n+1}$ and $r_n-u\in C_{n+1}$. For that candidate,
\[
 s\bigl(u+(r_n-u)\bigr)=r_n\in C_{n+2},
\]
because $r_n$ is prime. This contradicts its definition and proves
$d=c$.
\end{proof}

\section{Cardinalities, limiting density and proof of Theorem~\ref{thm:main}}

\begin{proof}[Completion of the proof of Theorem~\ref{thm:main}]
Propositions~\ref{prop:constants} and \ref{prop:equal} give
$M_n\sim Q_n\sim c\ph^n$. Since both $M_{n-1}$ and $Q_{n-1}$
are integers, Lemma~\ref{lem:filling} gives
\begin{align*}
 H_n(M_{n-1})
 &\le M_{n-1}-Q_{n-1}+H_n(Q_{n-1})\\
 &\le M_{n-1}-Q_{n-1}
       +O\left(\frac{Q_{n-1}}{\log^4Q_{n-1}}\right)
 =o(M_{n-1}).
\end{align*}
On the other hand, \eqref{eq:structure} implies
\[
 M_{n-1}-H_n(M_{n-1})\le |C_n|\le M_{n-1}+\pi(M_n).
\]
Because $M_n/M_{n-1}\to\ph$ and
$\pi(M_n)=O(M_n/\log M_n)=o(M_{n-1})$, these two bounds give
$|C_n|\sim M_{n-1}\sim c\ph^{n-1}$. Taking successive ratios and
dividing by $M_n$ proves \eqref{eq:main-limits}.
\end{proof}

\begin{remark}
The assertion $Q_n/M_n\to1$ says that the \emph{complete} prime
prefix reaches a relative distance $o(1)$ from the maximum. For the
integers we prove $H_n(M_{n-1})=o(M_{n-1})$, allowing a sparse set
of holes. In particular, the proof does not require the solid core
$K_n=\max\{K:\{1,\ldots,K\}\subseteq C_n\}$ to satisfy
$K_n\sim M_{n-1}$.
\end{remark}

\appendix
\section{Proof of the interval-restricted binary lemma}\label{app:binary}

We prove Lemma~\ref{lem:binary}, retaining its fixed parameters.
Only the two analytic estimates stated next are used without proof.
All subsequent interval, local-factor, and exceptional-set calculations
are included. Write $e(z)=\exp(2\pi iz)$, and use $\tot$ for Euler's
totient, to distinguish it from the golden ratio $\ph$.

\subsection{The analytic estimates and representation count}

The major-arc estimate we use is \cite[Lemma 5 (Major arc estimates)]{LyallRice}  or \cite[Proposition~24]{Tao}:
for an interval $\mathcal I\subseteq[1,x]$, a reduced fraction $a/q$,
and real $\xi$, one has, for every $D>0$,
\begin{equation}\label{eq:tao-input}
 \sum_{m\in\mathcal I\cap\Z}\Lambda(m)
       \ee{(a/q+\xi/q)m}
 =\frac{\mu(q)}{\tot(q)}\int_{\mathcal I}\ee{\xi v/q}\,dv
 +O_D\left((q+x|\xi|)\frac{x}{\log^D x}\right).
\end{equation}
Its implicit constant need not be effective. Here $\mu$ is the
M\"obius function and $\Lambda$ is the von Mangoldt function.

The minor-arc estimate is the following consequence of
\cite[Lemma~1]{Kumchev}: if $(a,q)=1$, $1\le q\le y$, and
$|\theta-a/q|\le q^{-2}$, then
\begin{equation}\label{eq:vaughan-input}
 \left|\sum_{p\le y}(\log p)e(p\theta)\right|
 \ll\bigl(yq^{-1/2}+y^{4/5}+(yq)^{1/2}\bigr)(\log(2y))^4.
\end{equation}
The additional factor $1+q^2|\theta-a/q|$ in the cited lemma is
at most $2$ under this hypothesis.

Define
\[
 S_I(\theta)=\sum_{p\in\Pp\cap I_X}(\log p)e(p\theta),\qquad
 S_J(\theta)=\sum_{q\in\Pp\cap J_X}(\log q)e(q\theta).
\]
Orthogonality of the functions $e(m\theta)$ for integers $m$ gives
\begin{equation}\label{eq:weighted-count}
 R_X(N):=\int_0^1 S_I(2\theta)S_J(\theta)e(-N\theta)\,d\theta
 =\sum_{\substack{p\in\Pp\cap I_X, q\in\Pp\cap J_X\\2p+q=N}}
       (\log p)(\log q).
\end{equation}
Thus $R_X(N)>0$ is equivalent to a representation of the required kind.
Fix
\begin{equation}\label{eq:choose-P}
 K=A+20,\qquad P=(\log X)^K.
\end{equation}
For reduced fractions $a/q$ with $1\le q\le P$ and $0\le a<q$,
let
\[
 \mathfrak M(a,q)=
 \{\theta\in\R/\Z:\|\theta-a/q\|\le P/X\},
 \qquad
 \mathfrak M=\bigcup_{q\le P}\ \bigcup_{\substack{0\le a<q\\(a,q)=1}}
                 \mathfrak M(a,q).
\]
Here $\|\cdot\|$ is circular distance to $0$; the fraction $0/1$
is included. Distinct fractions are at circular distance at least
$P^{-2}$, so these arcs are disjoint once $2P^3<X$.
Their total measure is $O(P^3/X)$. Set
$\mathfrak m=(\R/\Z)\setminus\mathfrak M$, and denote the two
parts of \eqref{eq:weighted-count} by $R_{\mathfrak M}(N)$ and
$R_{\mathfrak m}(N)$.

\subsection{The minor arcs}

Put $D_0=\ceil{X/P}$. Dirichlet approximation gives, for every
$\theta\in\mathfrak m$, a reduced fraction $a/q$ with
\[
 1\le q\le D_0,\qquad
 |\theta-a/q|\le\frac1{qD_0}\le\frac{P}{qX}.
\]
Distances may be interpreted modulo $1$ by changing $a$ by a
multiple of $q$. If $q\le P$, this places $\theta$ in
$\mathfrak M(a,q)$, a contradiction. Thus, for large $X$,
\begin{equation}\label{eq:minor-denominator}
 P<q\le\frac{2X}{P},\qquad |\theta-a/q|\le q^{-2}.
\end{equation}
Apply \eqref{eq:vaughan-input} to the two initial prime sums defining
$S_J$. Their endpoints are fixed positive multiples of $X$, so
$q$ is below both endpoints for all sufficiently large $X$.
An endpoint convention changes the sum by at most $O(\log X)$.
Using \eqref{eq:minor-denominator}, we obtain
\begin{equation}\label{eq:minor-sup}
 \sup_{\theta\in\mathfrak m}|S_J(\theta)|
 \ll XP^{-1/2}(\log X)^4+X^{4/5}(\log X)^4.
\end{equation}
Parseval's identity and \eqref{eq:theta} give
\[
 \int_0^1|S_I(2\theta)|^2\,d\theta
 =\sum_{p\in\Pp\cap I_X}(\log p)^2\ll X\log X.
\]
The factor $2$ does not affect orthogonality, since $2(p-p')$ is
a nonzero integer when $p\ne p'$. Bessel's inequality, applied to
$\ind_{\mathfrak m}(\theta)S_I(2\theta)S_J(\theta)$, therefore yields
\begin{align}
 \sum_{N\in\Z}|R_{\mathfrak m}(N)|^2
 &\le\int_{\mathfrak m}|S_I(2\theta)S_J(\theta)|^2\,d\theta\notag\\
 &\ll X^3P^{-1}(\log X)^9+X^{13/5}(\log X)^9
 \ll_A \frac{X^3}{\log^{A+2}X}.\label{eq:minor-L2}
\end{align}
For the last inequality, the first term has logarithmic exponent
$9-K=-A-11$, and the second satisfies
$X^{-2/5}\log^{A+11}X\to0$ after division by $X^3/\log^{A+2}X$.
In particular, for each fixed $\lambda>0$,
\begin{equation}\label{eq:minor-exceptions}
 \#\{N\in\Z:|R_{\mathfrak m}(N)|>\lambda X\}
 \ll_{A,\lambda}\frac{X}{\log^{A+2}X}.
\end{equation}

\subsection{The major arcs and the real solution measure}

For real $\zeta$, put
\[
 V_I(\zeta)=\int_{I_X}e(v\zeta)\,dv,\qquad
 V_J(\zeta)=\int_{J_X}e(v\zeta)\,dv.
\]
On an arc $\theta=a/q+\zeta$, $|\zeta|\le P/X$, let
$q'=q/(q,2)$, the denominator after reducing $2a/q$.
Formula \eqref{eq:tao-input} implies, for every fixed $E>0$,
\begin{equation}\label{eq:major-approx}
 \begin{aligned}
 S_J(a/q+\zeta)
  &=\frac{\mu(q)}{\tot(q)}V_J(\zeta)+O_{E,K}(X/\log^E X),\\
 S_I(2a/q+2\zeta)
  &=\frac{\mu(q')}{\tot(q')}V_I(2\zeta)+O_{E,K}(X/\log^E X).
 \end{aligned}
\end{equation}
Here is the required uniformity check. Choose $x=CX$ with fixed
$C>\max(\beta,\delta,1)$. In the first line of \eqref{eq:major-approx},
use $\xi=q\zeta$ in \eqref{eq:tao-input}; in the second, use
$\xi=2q'\zeta$ with denominator $q'$. In both cases
$q+x|\xi|=O(P^2)$, with $q'$ replacing $q$ when appropriate.
Taking exponent $D=E+2K+2$ in \eqref{eq:tao-input} gives the
stated error. Removing prime powers costs at most
$O(\sqrt X\log^2X)=O_E(X/\log^E X)$, as in Lemma~\ref{lem:pnt}.

Take $E=A+3K+10$. The sums $S_I,S_J$ and the integrals $V_I,V_J$
are $O(X)$; for the sums, this follows from \eqref{eq:theta}.
Since the major arcs have measure $O(P^3/X)$, multiplying
\eqref{eq:major-approx} and integrating causes total error
\begin{equation}\label{eq:major-product-error}
 O\bigl(XP^3/\log^E X\bigr)=O(X/\log^{A+10}X).
\end{equation}
Define the Ramanujan sum and its coefficient by
\[
 c_q(N)=\sum_{\substack{0\le a<q\\(a,q)=1}}e(-Na/q),\qquad
 w_q=\frac{\mu(q)\mu(q')}{\tot(q)\tot(q')}.
\]
The main term is
\[
 \sum_{q\le P}w_qc_q(N)
 \int_{-P/X}^{P/X}V_I(2\zeta)V_J(\zeta)e(-N\zeta)\,d\zeta.
\]
For any interval $[a,b]$ and nonzero $\zeta$,
$|\int_a^b e(v\zeta)\,dv|\le\min(b-a,1/(\pi|\zeta|))$.
Consequently
\[
 |V_I(2\zeta)V_J(\zeta)|\ll\min(X^2,|\zeta|^{-2}),
 \qquad
 \int_{|\zeta|>P/X}|V_I(2\zeta)V_J(\zeta)|\,d\zeta\ll X/P.
\]
Using $|c_q(N)|\le\tot(q)$, the error from extending every integral
to $\R$ is at most
\begin{equation}\label{eq:integral-tail}
 \frac{CX}{P}\sum_{q\le P}\frac1{\tot(q')}
 \ll\frac{X\log^2(2P)}{P}.
\end{equation}
Indeed each integer $q'$ arises from at most two values of $q$,
and $r/\tot(r)\le\tau(r)$, where $\tau$ counts divisors. Hence
\[
 \sum_{q\le P}\frac1{\tot(q')}
 \le2\sum_{r\le P}\frac{\tau(r)}r
 =2\sum_{ab\le P}\frac1{ab}
 \le2(1+\log P)^2.
\]
The inequality $r/\tot(r)\le\tau(r)$ follows prime by prime:
for $\ell^a\parallel r$, $\ell/(\ell-1)\le2\le a+1$.

For clarity, Fourier inversion introduces no unspecified geometric
constant here. The continuous function
\[
 f(y)=\left(\tfrac12\ind_{2I_X}*\ind_{J_X}\right)(y)
 =\int_{I_X}\ind_{J_X}(y-2v)\,dv
\]
has Fourier transform $V_I(2\zeta)V_J(\zeta)$ with the positive
exponential convention used above. That product is integrable by
the displayed bound, and $f(N)=\mathcal J_X(N)$. Fourier inversion gives
\begin{equation}\label{eq:singular-integral}
 \int_\R V_I(2\zeta)V_J(\zeta)e(-N\zeta)\,d\zeta
 =\mathcal J_X(N).
\end{equation}
Write
\[
 \mathfrak S_P(N)=\sum_{q\le P}w_qc_q(N).
\]
Combining \eqref{eq:major-product-error}--\eqref{eq:singular-integral},
and using $K=A+20$, proves the uniform formula
\begin{equation}\label{eq:major-final}
 R_{\mathfrak M}(N)
 =\mathfrak S_P(N)\mathcal J_X(N)+O_A(X/\log^{A+2}X).
\end{equation}
For the second error, note explicitly that
$\log^2(2P)/P=O_A(\log^{-A-2}X)$.

\subsection{An elementary bound for the series tail}

We first prove that, for $Y\ge2$ and $N\ge1$,
\begin{equation}\label{eq:series-tail}
 \sum_{q>Y}\frac{|c_q(N)|}{\tot(q)^2}
 \ll\frac{\log^3(2Y)}{Y}
       F(N),\qquad
 F(N)=\sum_{h\mid N}\left(\frac{h}{\tot(h)}\right)^2.
\end{equation}
The identity
\begin{equation}\label{eq:ramanujan-identity}
 c_q(N)=\sum_{h\mid(q,N)}h\mu(q/h)
\end{equation}
follows by inserting
$\ind_{(a,q)=1}=\sum_{d\mid(a,q)}\mu(d)$ into the defining exponential
sum. The sum over multiples of $d$ is $q/d$ if $q/d\mid N$ and
zero otherwise; set $h=q/d$. In particular,
$|c_q(N)|\le\sum_{h\mid(q,N)}h$.

The product formula for the totient gives
$\tot(hk)\ge\tot(h)\tot(k)$. Therefore, writing $q=hk$ and
interchanging nonnegative sums,
\begin{equation}\label{eq:tail-divisors}
 \sum_{q>Y}\frac{|c_q(N)|}{\tot(q)^2}
 \le\sum_{h\mid N}\frac{h}{\tot(h)^2}
             \sum_{k>Y/h}\frac1{\tot(k)^2}.
\end{equation}
For $z\ge1$ we claim
\begin{equation}\label{eq:totient-tail}
 \sum_{k>z}\frac1{\tot(k)^2}\ll\frac{\log^3(2z)}z.
\end{equation}
Let $\tau_4(k)$ count ordered factorizations $k=abcd$ in positive
integers. Prime by prime,
\[
 \left(\frac{k}{\tot(k)}\right)^2\le\tau(k)^2\le\tau_4(k),
\]
since $(a+1)^2\le\binom{a+3}{3}$ for every integer $a\ge0$.
Furthermore, for $w\ge1$,
\[
 \sum_{k\le w}\tau_4(k)
 =\sum_{abc\le w}\floor{\frac{w}{abc}}
 \le w\left(\sum_{a\le w}\frac1a\right)^3
 \le w(1+\log w)^3.
\]
Split $k>z$ into ranges $2^jz<k\le2^{j+1}z$. These inequalities give
\[
 \sum_{k>z}\frac1{\tot(k)^2}
 \le\sum_{j\ge0}\frac1{(2^jz)^2}
                    \sum_{k\le2^{j+1}z}\tau_4(k)
 \ll\frac1z\sum_{j\ge0}2^{-j}(\log(2z)+j)^3
 \ll\frac{\log^3(2z)}z,
\]
because $\sum_{j\ge0}2^{-j}(j+1)^3$ converges. This proves
\eqref{eq:totient-tail}. For $0\le z<1$, the same sum is $O(1)$,
by the bound at $z=1$ and the single term $k=1$.

For $h\le Y$, apply \eqref{eq:totient-tail} to $z=Y/h$ in
\eqref{eq:tail-divisors}, bounding $\log(2Y/h)$ by $\log(2Y)$.
The resulting contribution for $h$ is at most
\[
 C\frac{\log^3(2Y)}Y\frac{h^2}{\tot(h)^2}.
\]
For $h>Y$, the $O(1)$ bound gives at most $Ch/\tot(h)^2$,
which is also bounded by the same expression after increasing $C$.
This proves \eqref{eq:series-tail}, including absolute convergence
for each fixed $N$.

We will also need its first moment. For large $X$,
\begin{align}
 \sum_{1\le N\le BX}F(N)
 &\le BX\sum_{h\le BX}\frac1h
                  \left(\frac{h}{\tot(h)}\right)^2\notag\\
 &\le BX\sum_{h\le BX}\frac{\tau_4(h)}h
 =BX\sum_{abcd\le BX}\frac1{abcd}
 \le BX(1+\log(BX))^4
 \ll_B X\log^4X.\label{eq:F-mean}
\end{align}

\subsection{The local factors and positivity}

For odd $N$, absolute convergence from \eqref{eq:series-tail}
allows us to compute the full series
$\mathfrak S(N)=\sum_{q\ge1}w_qc_q(N)$. A nonzero coefficient
requires $q$ squarefree. For odd squarefree $q$, we have $q'=q$
and $w_q=\mu(q)^2/\tot(q)^2$. For even squarefree $q=2r$,
with $r$ odd, we have $q'=r$ and
\[
 w_{2r}=-\frac{\mu(r)^2}{\tot(r)^2},\qquad
 c_{2r}(N)=c_2(N)c_r(N)=-c_r(N).
\]
Here $c_q(N)$ is multiplicative in $q$ by
\eqref{eq:ramanujan-identity}, and $c_2(N)=-1$ for odd $N$.
Thus the even term equals the corresponding odd term, and
\begin{equation}\label{eq:full-series}
 \begin{aligned}
 \mathfrak S(N)
 &=2\sum_{\substack{r\ge1\\r\text{ odd}}}
               \frac{\mu(r)^2c_r(N)}{\tot(r)^2}\\
 &=2\prod_{\ell>2}
             \left(1-\frac1{(\ell-1)^2}\right)
       \prod_{\substack{\ell\mid N\\\ell>2}}
             \frac{\ell-1}{\ell-2}.
 \end{aligned}
\end{equation}
All products indexed by $\ell$ here are over primes. To verify
each factor, \eqref{eq:ramanujan-identity} gives
$c_\ell(N)=\ell-1$ if $\ell\mid N$ and $c_\ell(N)=-1$ otherwise.
The factor at a divisor prime is $1+1/(\ell-1)$; dividing this by
$1-1/(\ell-1)^2$ gives $(\ell-1)/(\ell-2)$, as displayed.

In particular,
\begin{equation}\label{eq:s0}
 \mathfrak S(N)\ge s_0:=
 2\prod_{\ell>2}\left(1-\frac1{(\ell-1)^2}\right)>0
 \qquad(N\text{ odd}).
\end{equation}
For positivity of this constant, the factors lie in $(0,1)$,
$1/(\ell-1)^2\le1/4$, and
$\sum_{\ell>2}1/(\ell-1)^2<\infty$. The inequality
$\log(1-z)\ge-2z$ for $0\le z\le1/2$ bounds the product away
from zero.

The odd terms of $\mathfrak S_P(N)$ have $r\le P$, whereas its
even terms have $r\le P/2$. Consequently, for odd $N$,
\[
 |\mathfrak S(N)-\mathfrak S_P(N)|
 \le2\sum_{r>P/2}\frac{|c_r(N)|}{\tot(r)^2}
 \ll\frac{\log^3(2P)}P F(N).
\]
Summing over $1\le N\le BX$ and using \eqref{eq:F-mean}, we see
that the number of odd $N$ in this range for which
$|\mathfrak S(N)-\mathfrak S_P(N)|>s_0/2$ is at most
\begin{equation}\label{eq:series-exceptions}
 O_B\left(\frac{X\log^4X\log^3(2P)}P\right)
 =O_A\left(\frac{X}{\log^A X}\right).
\end{equation}
The last estimate follows from $K=A+20$, since
$\log^3(2P)=O_A((\log\log X)^3)$.

\subsection{Completion of the representation lemma}

Outside the exceptional set in \eqref{eq:series-exceptions}, every
odd target has $\mathfrak S_P(N)\ge s_0/2$. If also
$\mathcal J_X(N)\ge\eta X$, then \eqref{eq:major-final} gives,
for sufficiently large $X$,
\[
 \operatorname{Re}R_{\mathfrak M}(N)
 \ge\frac{s_0\eta X}{2}-\frac{s_0\eta X}{8}
 =\frac{3s_0\eta X}{8}.
\]
By \eqref{eq:minor-exceptions} with $\lambda=s_0\eta/8$,
all but $O_A(X/\log^{A+2}X)$ further targets satisfy
$|R_{\mathfrak m}(N)|\le s_0\eta X/8$. For a target outside
both exceptional sets, \eqref{eq:weighted-count} is a nonnegative
real number and obeys
\[
 R_X(N)=\operatorname{Re}
             \bigl(R_{\mathfrak M}(N)+R_{\mathfrak m}(N)\bigr)
 \ge\frac{s_0\eta X}{4}>0.
\]
The union of the exceptional sets has size $O_A(X/\log^A X)$.
This proves Lemma~\ref{lem:binary}. All interval data were fixed
before $X$ tended to infinity; no uniformity as an interval length
or margin tends to zero has been used. \hfill$\square$

\section*{Declaration of Generative AI and AI-assisted technologies in the writing process}
During the preparation of this work the author used GPT-6 Astra in order to assist with the construction of the mathematical argument and the drafting of the manuscript. After using this tool, the author reviewed and edited the content as needed and takes full responsibility for the content of the publication. The main theorem has been independently checked in Lean~4, as recorded in \cite{Github}.

\section*{Declarations of interest}
Declarations of interest: none.

\begin{thebibliography}{99}

\bibitem{BakerHarman1998}
R.~C. Baker and G. Harman,
\emph{The three primes theorem with almost equal summands},
Philos. Trans. Roy. Soc. London Ser. A
\textbf{356} (1998), no.~1738, 763--780.

\bibitem{CVZ}
M. Caragiu, P. A. Vicol, and M. Zaki,
\href{https://www.fq.math.ca/Papers1/55-4/CaragiuVicolZaki03162017.pdf}
{On Conway's subprime function, a covering of $\N$ and an unexpected
appearance of the golden ratio},
\emph{The Fibonacci Quarterly} \textbf{55} (2017), no.~4, 327--331.
%See Theorem~1 and Conjecture~3.

\bibitem{CL}
G. Coppola and M. B. S. Laporta,
%\href{https://www.giovannicoppola.name/files/articoli/8_giovanni_coppola_number_theory_seminar_politecnico_torino_245.pdf}
{On the representation of even integers as sum of two almost equal primes},
\emph{Rendiconti del Seminario Matematico, Universit\`a e Politecnico
di Torino} \textbf{53} (1995), no.~3, 245--252.
%See the corollary to Theorem~1, p.~246.

\bibitem{Github}
R.~Popescu,
\emph{Lean formalization of golden-ratio growth of Conway's subprime closure},
GitHub repository, 2026.
\url{https://github.com/Renriviera/conway-subprime-golden}.

\bibitem{GuyKhovanovaSalazar2014}
R.~K. Guy, T. Khovanova, and J. Salazar,
\emph{Conway's Subprime Fibonacci Sequences},
Mathematics Magazine \textbf{87} (2014), no.~5, 323--337.

\bibitem{Kumchev}
A. V. Kumchev,
\href{https://tigerweb.towson.edu/akumchev/a23.pdf}
{On sums of primes from Beatty sequences},
\emph{Integers} \textbf{8} (2008), Article A08, 12 pp.
%See Lemma~1.


\bibitem{LyallRice}
N. Lyall and A. Rice, 
{Polynomial Differences in the Primes}, in \emph{Combinatorial and Additive Number Theory: CANT 2011 and 2012}, 
Springer Proceedings in Mathematics \& Statistics \textbf{101} (2014), 129--146.


\bibitem{OEIS}
OEIS Foundation Inc. (2011), The On-Line Encyclopedia of Integer Sequences, A117818, \href{https://oeis.org/A117818}.

\bibitem{openai2026astra}
OpenAI, \emph{GPT-6 Astra},
2026. \url{https://openai.com/index/gpt-6-astra/}.

\bibitem{Roberts}
S. Roberts, 
{Genius At Play: The Curious Mind of John Horton Conway, Audiobook, Jennifer Van Dyck
(Reader)}, 
Bloomsbury Press, 2015.

\bibitem{Tao}
T. Tao,
\href{https://terrytao.wordpress.com/2015/03/30/254a-notes-8-the-hardy-littlewood-circle-method-and-vinogradovs-theorem/}
{254A, Notes 8: The Hardy--Littlewood circle method and Vinogradov's theorem},
lecture notes, 30 March 2015.
%See Proposition~24. Accessed 8 September 2026.

\end{thebibliography}

\bigskip
\noindent
\begin{scriptsize}
\textsc{Department of Computer Science, Columbia University,
New York, NY 10027, USA}
%

\noindent
\textit{Email address:}  rep2159@columbia.edu
\end{scriptsize}
\end{document}
